\chapter{Reactiemechanismen: gebogen pijlen}\label{ch:g12:curly-arrows}

Een \omterm{def:g8:balanced-equations:equation}{kloppende reactievergelijking} is als de eerste en de laatste foto van
een film: ze laat de \omterm{def:g7:chemical-reactions:reactant}{beginstoffen} ervoor en de \omterm{def:g7:chemical-reactions:reactant}{reactieproducten} erna zien,
en niets daartussen. Toch ruilen \omterm{def:g7:atoms-and-molecules:molecule}{moleculen} niet zomaar \omterm{def:g7:atoms-and-molecules:atom}{atomen}. Bindingen
worden een voor een verbroken en gevormd, terwijl elektronenparen van de
ene plek naar de andere verschuiven, via kortlevende deeltjes die nooit
in de \omterm{def:g8:balanced-equations:equation}{reactievergelijking} verschijnen. Scheikundigen tekenen die
verschuivingen met \omterm{def:g12:curly-arrows:curly-arrow}{gebogen pijlen}. Met een paar regels verklaren die
pijlen waarom een reactie plaatsvindt waar ze plaatsvindt, en voorspellen
ze reacties die nog nooit gezien zijn.

\begin{recall}
\omterm{def:g11:polarity-and-cohesion:electronegativity}{Elektronegativiteit}, \omterm{def:g11:polarity-and-cohesion:polar-bond}{polaire bindingen} en \omterm{def:g11:polarity-and-cohesion:polar-bond}{partiële ladingen}
(\cref{ch:g11:polarity-and-cohesion}). \omterm{def:g10:lewis-and-shape:lewis-structure}{Lewisstructuren}: bindende en \omterm{def:g10:lewis-and-shape:lone-pair}{vrije
elektronenparen} (\cref{ch:g10:lewis-and-shape}). \omterm{def:g11:functional-groups:functional-group}{Karakteristieke groepen}
en families (\cref{ch:g11:functional-groups}).
\end{recall}

\section{Elektronenrijke en elektronenarme plaatsen}

\begin{definition}[Donorplaatsen en acceptorplaatsen]\label{def:g12:curly-arrows:site}
Een \emph{donorplaats}\index{donorplaats} van een \omterm{def:g7:atoms-and-molecules:molecule}{molecuul} of een \omterm{def:g9:ions:ion}{ion} is
een plek die rijk is aan \omterm{def:g9:inside-the-atom:nucleus}{elektronen} en een paar kan afstaan om een nieuwe
binding te vormen: een \omterm{def:g10:lewis-and-shape:lone-pair}{vrij elektronenpaar}, een \omterm{def:g10:lewis-and-shape:multiple-bond}{dubbele binding}, een \omterm{def:g7:atoms-and-molecules:atom}{atoom}
met een negatieve lading of een \omterm{def:g11:polarity-and-cohesion:polar-bond}{partiële lading} $\delta^-$. Een
\emph{acceptorplaats}\index{acceptorplaats} is een plek die arm is aan
\omterm{def:g9:inside-the-atom:nucleus}{elektronen} en een paar kan opnemen: een \omterm{def:g7:atoms-and-molecules:atom}{atoom} met een positieve lading of
een \omterm{def:g11:polarity-and-cohesion:polar-bond}{partiële lading}
$\delta^+$.
\end{definition}

\begin{example}[Plaatsen in twee moleculen]\label{ex:g12:curly-arrows:sites}
In broomethaan, \ce{CH3-CH2-Br}, is broom elektronegatiever dan koolstof:
het koolstofatoom dat eraan gebonden is, draagt $\delta^+$
(\omterm{def:g12:curly-arrows:site}{acceptorplaats}), het broomatoom $\delta^-$ en drie \omterm{def:g10:lewis-and-shape:lone-pair}{vrije elektronenparen}.
In propanon draagt het zuurstofatoom van de \ce{C=O}-groep $\delta^-$ en
twee \omterm{def:g10:lewis-and-shape:lone-pair}{vrije elektronenparen} (donorplaatsen), zijn koolstofatoom $\delta^+$
(\omterm{def:g12:curly-arrows:site}{acceptorplaats}). Het \omterm{def:g9:ions:ion}{hydroxide-ion} \ce{OH-}, met zijn negatieve lading en
drie \omterm{def:g10:lewis-and-shape:lone-pair}{vrije elektronenparen}, is een sterke donor.
\end{example}

\begin{omfigure}
\setchemfig{atom sep=2.2em}
\begin{tabular}{cc}
\chemfig{H_3C-C(-[2]H)(-[6]H)-\charge{90:2pt=\:,0=\:,270:2pt=\:}{Br}} &
\chemfig{H_3C-C(-[6]CH_3)=[2]\charge{0=\:,180=\:}{O}} \\[2pt]
\small $\delta^+$ op het koolstofatoom gebonden aan \ce{Br}, $\delta^-$ op \ce{Br} &
\small $\delta^+$ op het koolstofatoom van \ce{C=O}, $\delta^-$ op \ce{O} \\
\small broomethaan & \small propanon
\end{tabular}

{\small Elektronenarme koolstofatomen ($\delta^+$, acceptorplaatsen) naast
elektronenrijke, elektronegatieve \omterm{def:g7:atoms-and-molecules:atom}{atomen} ($\delta^-$, met \omterm{def:g10:lewis-and-shape:lone-pair}{vrije
elektronenparen}: donorplaatsen).}
\end{omfigure}

\section{Gebogen pijlen}

\begin{definition}[Gebogen pijl]\label{def:g12:curly-arrows:curly-arrow}
Een \emph{gebogen pijl}\index{gebogen pijl} laat de verplaatsing zien van
een elektronenpaar tijdens een stap van een reactie. Hij begint bij het
paar dat zich verplaatst, een \omterm{def:g10:lewis-and-shape:lone-pair}{vrij elektronenpaar} of een binding, en
eindigt waar het paar naartoe gaat: op een \omterm{def:g7:atoms-and-molecules:atom}{atoom}, waar het een nieuwe
binding vormt, of op een \omterm{def:g7:atoms-and-molecules:atom}{atoom} dat het paar overneemt van een binding die
verbroken wordt.
\end{definition}

\begin{method}[Gebogen pijlen tekenen]\label{met:g12:curly-arrows:drawing}
\begin{enumerate}
  \item Zoek de \omterm{def:g12:curly-arrows:site}{donorplaats} en de \omterm{def:g12:curly-arrows:site}{acceptorplaats}.
  \item Teken een pijl van het donorpaar (een \omterm{def:g10:lewis-and-shape:lone-pair}{vrij elektronenpaar} of een
        binding) naar het acceptoratoom: daar ontstaat een nieuwe binding.
  \item Als het acceptoratoom dan te veel \omterm{def:g9:inside-the-atom:nucleus}{elektronen} zou krijgen (meer
        dan een octet, of meer dan twee bij waterstof), teken dan een
        tweede pijl die een van zijn bindingen verbreekt, waarbij het
        paar naar het elektronegatiefste \omterm{def:g7:atoms-and-molecules:atom}{atoom} gaat.
  \item Controleer de ladingen na de stap: het \omterm{def:g7:atoms-and-molecules:atom}{atoom} dat een \omterm{def:g10:lewis-and-shape:lone-pair}{vrij
        elektronenpaar} heeft afgestaan, wordt één eenheid positiever, het
        \omterm{def:g7:atoms-and-molecules:atom}{atoom} dat een paar als \omterm{def:g10:lewis-and-shape:lone-pair}{vrij elektronenpaar} heeft opgenomen één
        eenheid negatiever; de totale lading verandert niet.
\end{enumerate}
\end{method}

\begin{proposition}[In elke stap blijft de lading behouden]\label{prop:g12:curly-arrows:charge}
De totale elektrische lading van de deeltjes is voor en na elke stap van
een mechanisme dezelfde.
\end{proposition}

\begin{proof}
Een \omterm{def:g12:curly-arrows:curly-arrow}{gebogen pijl} verplaatst alleen \omterm{def:g9:inside-the-atom:nucleus}{elektronen} die er al zijn; er wordt
geen \omterm{def:g9:inside-the-atom:nucleus}{elektron} gemaakt of vernietigd, en de kernen veranderen niet.
\end{proof}

\begin{example}[Een eerste pijl]\label{ex:g12:curly-arrows:first}
Een waterstofion \ce{H+} ontmoet een \omterm{def:g9:ions:ion}{hydroxide-ion} \ce{OH-}. Een \omterm{def:g10:lewis-and-shape:lone-pair}{vrij
elektronenpaar} van het zuurstofatoom vormt een binding met het
waterstofion, dat geen \omterm{def:g9:inside-the-atom:nucleus}{elektronen} heeft:
\ce{H+ + OH- -> H2O}. Eén pijl, van het \omterm{def:g10:lewis-and-shape:lone-pair}{vrije elektronenpaar} van het
zuurstofatoom naar het \ce{H+}; de ladingen $+1$ en $-1$ worden 0 in het
watermolecuul.
\end{example}

\begin{omfigure}
\setchemfig{atom sep=2.1em}
\chemfig{@{hp}H^{+}}\qquad\raisebox{0.2ex}{$+$}\qquad
\chemfig{H-@{ox}\charge{90=\:,0=\:,270=\:}{O}^{-}}
\qquad\raisebox{0.2ex}{$\longrightarrow$}\qquad
\chemfig{H-\charge{90=\:,270=\:}{O}-H}
\chemmove{
  \draw[->, thick, omProp, shorten <=4pt, shorten >=3pt]
    ($(ox)+(0,0.25)$) .. controls +(110:9mm) and +(60:9mm) .. (hp);}

{\small De eenvoudigste stap: een \omterm{def:g10:lewis-and-shape:lone-pair}{vrij elektronenpaar} van het
\omterm{def:g9:ions:ion}{hydroxide-ion} verplaatst zich en vormt de nieuwe \ce{O-H}-binding met het
waterstofion.}
\end{omfigure}

\section{Elementaire stappen en intermediairen}

\begin{definition}[Mechanisme, elementaire stap, intermediair]\label{def:g12:curly-arrows:mechanism}
Het \emph{reactiemechanisme}\index{reactiemechanisme} van een reactie is
de reeks stappen waarlangs de \omterm{def:g7:chemical-reactions:reactant}{beginstoffen} de \omterm{def:g7:chemical-reactions:reactant}{reactieproducten} worden.
Elke \emph{elementaire stap}\index{elementaire stap} is één gebeurtenis
waarbij een paar elektronenparen tegelijk verschuiven. Een deeltje dat in
de ene stap ontstaat en in een latere stap verbruikt wordt, is een
\emph{reactie-intermediair}\index{reactie-intermediair}:
het komt niet voor in de totale \omterm{def:g8:balanced-equations:equation}{reactievergelijking}.
\end{definition}


\begin{example}[Een carbokation als intermediair]\label{ex:g12:curly-arrows:carbocation}
Als waterstofbromide aan etheen addeert, verloopt de reactie in twee
stappen. In de eerste geeft de \omterm{def:g10:lewis-and-shape:multiple-bond}{dubbele binding} van etheen een paar af aan
het waterstofatoom van \ce{HBr}, en het paar van \ce{H-Br} gaat naar het
broomatoom, dat vertrekt als \ce{Br-}; het koolstofatoom dat zijn aandeel
in de \omterm{def:g10:lewis-and-shape:multiple-bond}{dubbele binding} kwijt is, houdt maar zes \omterm{def:g9:inside-the-atom:nucleus}{elektronen} over en krijgt
een positieve lading: een carbokation,
\ce{CH3-CH2+}. In de tweede stap vormt een \omterm{def:g10:lewis-and-shape:lone-pair}{vrij elektronenpaar} van
\ce{Br-} een binding met dat koolstofatoom. Het carbokation is een
intermediair.
\end{example}

\section{Drie soorten reacties}

\begin{definition}[Substitutie, additie, eliminatie]\label{def:g12:curly-arrows:categories}
Bij een \emph{substitutie}\index{substitutie} wordt een \omterm{def:g7:atoms-and-molecules:atom}{atoom} of groep
van een \omterm{def:g7:atoms-and-molecules:molecule}{molecuul} door een ander vervangen. Bij een
\emph{additie}\index{additie} worden twee \omterm{def:g7:atoms-and-molecules:molecule}{moleculen} één, waarbij een
meervoudige binding een enkelvoudige wordt. Bij een
\emph{eliminatie}\index{eliminatie} wordt een klein \omterm{def:g7:atoms-and-molecules:molecule}{molecuul} (vaak water)
van een \omterm{def:g7:atoms-and-molecules:molecule}{molecuul} afgesplitst, waarbij een enkelvoudige binding een
meervoudige wordt.
\end{definition}

\begin{omfigure}
\setchemfig{atom sep=2em}
\small
\textbf{\omterm{def:g12:curly-arrows:categories}{Substitutie}} (één stap): \ce{CH3-CH2-Br + OH- -> CH3-CH2-OH + Br-}

\medskip
\chemfig{H-@{o1}\charge{90=\:,0=\:,270=\:}{O}^{-}}\qquad
\chemfig{H_3C-@{c1}C(-[2]H)(-[6]H)-[@{b1}]@{br1}\charge{90=\:,0=\:,270=\:}{Br}}
\qquad$\longrightarrow$\qquad
\chemfig{H_3C-C(-[2]H)(-[6]H)-O-H}\quad$+$\quad\ce{Br-}
\chemmove{
  \draw[->, thick, omProp, shorten <=4pt, shorten >=3pt]
    ($(o1)+(0.15,0.2)$) .. controls +(60:9mm) and +(120:9mm) .. (c1);
  \draw[->, thick, omDef, shorten <=2pt, shorten >=4pt]
    (b1) .. controls +(-90:7mm) and +(-90:7mm) .. (br1);}

\vspace{6mm}
\textbf{\omterm{def:g12:curly-arrows:categories}{Additie}} (twee stappen): \ce{CH2=CH2 + HBr -> CH3-CH2-Br}

\vspace{9mm}
\chemfig{H_2C=[@{pi}]CH_2}\quad$+$\quad
\chemfig{@{h2}H-[@{b2}]@{br2}Br}
\qquad$\longrightarrow$\qquad
\chemfig{H_3C-@{cp}\charge{90:2pt=$\scriptstyle +$}{C}H_2}\quad$+$\quad
\chemfig{@{brm}\charge{90=\:,180=\:,270=\:,0=\:}{Br}^{-}}
\qquad$\longrightarrow$\qquad \ce{CH3-CH2-Br}
\chemmove{
  \draw[->, thick, omProp, shorten <=2pt, shorten >=3pt]
    (pi) .. controls +(90:10mm) and +(110:9mm) .. (h2);
  \draw[->, thick, omDef, shorten <=2pt, shorten >=4pt]
    (b2) .. controls +(-90:7mm) and +(-90:7mm) .. (br2);
  \draw[->, thick, omProp, shorten <=4pt, shorten >=3pt]
    ($(brm)+(0,-0.25)$) .. controls +(-120:8mm) and +(-60:8mm) .. (cp);}

\vspace{6mm}
\textbf{\omterm{def:g12:curly-arrows:categories}{Eliminatie}} (zure dehydratatie van ethanol, drie stappen):
\ce{CH3-CH2-OH -> CH2=CH2 + H2O}

\medskip
\begin{tabular}{@{}l@{}}
1.\ het zuurstofatoom neemt een waterstofion op: \ce{CH3-CH2-OH + H+ -> CH3-CH2-OH2+};\\
2.\ het \ce{C-O}-paar vertrekt met water: \ce{CH3-CH2-OH2+ -> CH3-CH2+ + H2O};\\
3.\ een naburig \ce{C-H}-paar wordt de \omterm{def:g10:lewis-and-shape:multiple-bond}{dubbele binding}, en het
    \ce{H+} vertrekt:\\
\phantom{3.\ }\ce{CH3-CH2+ -> CH2=CH2 + H+}.
\end{tabular}

\medskip
{\small Drie mechanismen. Oranje pijlen: een donorpaar vormt een nieuwe
binding. Blauwe pijlen: een binding wordt verbroken en haar paar gaat
naar het elektronegatiefste \omterm{def:g7:atoms-and-molecules:atom}{atoom}. Bij de \omterm{def:g12:curly-arrows:categories}{eliminatie} wordt het
waterstofion aan het eind teruggegeven.}
\end{omfigure}


\begin{remark}[Het skelet of de groep veranderen]\label{rem:g12:curly-arrows:skeleton}
Sommige reacties veranderen alleen de \omterm{def:g11:functional-groups:functional-group}{karakteristieke groep} en laten het
koolstofskelet intact: de \omterm{def:g12:curly-arrows:categories}{substitutie} van broomethaan maakt van een
\omterm{def:g11:functional-groups:halogenoalkane}{halogeenalkaan} een \omterm{def:g11:functional-groups:alcohol}{alcohol}, met twee koolstofatomen ervoor en twee erna.
Andere bouwen het skelet op of knippen het door: een nieuwe
koolstof-koolstofbinding maken verlengt een keten, de sleutelstap bij het
maken van een groot \omterm{def:g7:atoms-and-molecules:molecule}{molecuul} uit kleine. Een \omterm{def:g10:synthesis-yield:synthesis}{synthese} plannen betekent
reacties van beide soorten kiezen, in de juiste volgorde.
\end{remark}

\section{Oefeningen}

\begin{exercise}[$\star$]\label{exo:g12:curly-arrows:1}
Noem in elk deeltje een \omterm{def:g12:curly-arrows:site}{donorplaats} en, als die er is, een
\omterm{def:g12:curly-arrows:site}{acceptorplaats}: \ce{OH-}; \ce{H+}; \ce{CH2=CH2}; \ce{CH3-Cl}; \ce{H2O}.
\end{exercise}

\begin{exercise}[$\star$]\label{exo:g12:curly-arrows:2}
\omterm{def:g12:curly-arrows:categories}{Substitutie}, \omterm{def:g12:curly-arrows:categories}{additie} of \omterm{def:g12:curly-arrows:categories}{eliminatie}?
(a) \ce{CH3-CH2-Cl + OH- -> CH3-CH2-OH + Cl-};
(b) \ce{CH2=CH2 + H2O -> CH3-CH2-OH};
(c) \ce{CH3-CH2-OH -> CH2=CH2 + H2O};
(d) \ce{CH2=CH2 + Br2 -> CH2Br-CH2Br}.
\end{exercise}

\begin{exercise}[$\star$]\label{exo:g12:curly-arrows:3}
Waar begint een \omterm{def:g12:curly-arrows:curly-arrow}{gebogen pijl}, en waar eindigt hij? Wat betekent een
\omterm{def:g12:curly-arrows:curly-arrow}{gebogen pijl} van een \omterm{def:g10:lewis-and-shape:lone-pair}{vrij elektronenpaar} van een zuurstofatoom naar een
koolstofatoom?
\end{exercise}

\begin{exercise}[$\star$]\label{exo:g12:curly-arrows:4}
Welk deeltje geeft in de reactie \ce{H+ + NH3 -> NH4+} het paar af?
Teken de pijl. Wat is de lading van het product?
\end{exercise}

\begin{exercise}[$\star$]\label{exo:g12:curly-arrows:5}
Wat is een \omterm{def:g12:curly-arrows:mechanism}{reactie-intermediair}? Noem het intermediair van de \omterm{def:g12:curly-arrows:categories}{additie} van
\ce{HBr} aan etheen.
\end{exercise}

\begin{exercise}[$\star\star$]\label{exo:g12:curly-arrows:6}
Teken de \omterm{def:g12:curly-arrows:curly-arrow}{gebogen pijlen} van de \omterm{def:g12:curly-arrows:categories}{substitutie}
\ce{CH3-Br + OH- -> CH3-OH + Br-} (één stap), en controleer de ladingen.
\end{exercise}

\begin{exercise}[$\star\star$]\label{exo:g12:curly-arrows:7}
Waarom gaat in de eerste stap van de \omterm{def:g12:curly-arrows:categories}{additie} van \ce{HBr} aan etheen het
\ce{H-Br}-paar naar het broomatoom en niet naar het waterstofatoom?
Welke lading draagt elk product van die stap?
\end{exercise}

\begin{exercise}[$\star\star$]\label{exo:g12:curly-arrows:8}
Schrijf bij de dehydratatie van ethanol de drie stappen op en zeg bij
elke stap welk deeltje de donor is en welk de acceptor. Waarom staat het
waterstofion niet in de totale \omterm{def:g8:balanced-equations:equation}{reactievergelijking}?
\end{exercise}

\begin{exercise}[$\star\star$]\label{exo:g12:curly-arrows:9}
De stap \ce{CH3-CH2+ + H2O -> CH3-CH2-OH2+} volgt op de
\omterm{def:g12:curly-arrows:categories}{additie} van een waterstofion aan etheen in water. Teken de pijl. Wat
ontstaat er na het afsplitsen van nog een \ce{H+}? Schrijf de totale
\omterm{def:g8:balanced-equations:equation}{reactievergelijking} op en geef aan tot welke soort ze behoort.
\end{exercise}

\begin{exercise}[$\star\star$]\label{exo:g12:curly-arrows:10}
Lees het mechanisme van de \omterm{def:g12:curly-arrows:categories}{substitutie} van broomethaan: hoeveel paren
verschuiven er? Welke binding ontstaat, welke wordt verbroken? Is er een
intermediair?
\end{exercise}

\begin{exercise}[$\star\star$]\label{exo:g12:curly-arrows:11}
% ledger: en:C, en:Cl, en:Br
Chloormethaan reageert langzamer met het \omterm{def:g9:ions:ion}{hydroxide-ion} dan
broommethaan. Leg met de \omterm{def:g11:polarity-and-cohesion:electronegativity}{elektronegativiteiten} (C 2.55, Cl 3.16, Br
2.96) uit welk koolstofatoom het meest $\delta^+$ is. Verklaart dat de
snelheden? (Denk ook aan hoe gemakkelijk het vertrekkende \omterm{def:g9:ions:ion}{ion} het paar
meeneemt.)
\end{exercise}

\begin{exercise}[$\star\star\star$]\label{exo:g12:curly-arrows:12}
Waterstofchloride addeert aan propeen, \ce{CH2=CH-CH3}. In de eerste stap
kan het waterstofion aan elk van de twee koolstofatomen van de \omterm{def:g10:lewis-and-shape:multiple-bond}{dubbele
binding} binden, wat twee mogelijke carbokationen geeft. Teken ze allebei,
en de twee mogelijke producten, en geef hun namen. (Welk het meest
ontstaat, wordt bestudeerd in het deel voor het eerste
universiteitsjaar.)
\end{exercise}

\begin{exercise}[$\star\star\star$]\label{exo:g12:curly-arrows:13}
Een leerling tekent de eerste stap van de \omterm{def:g12:curly-arrows:categories}{substitutie} van broomethaan met
een pijl van het broomatoom naar het zuurstofatoom van \ce{OH-}. Leg de
twee fouten uit.
\end{exercise}

\begin{exercise}[$\star\star\star$]\label{exo:g12:curly-arrows:14}
Bij de vorming van een \omterm{def:g11:functional-groups:carboxylic-acid}{ester} uit een \omterm{def:g11:functional-groups:carboxylic-acid}{carbonzuur} en een \omterm{def:g11:functional-groups:alcohol}{alcohol} valt het
zuurstofatoom van de \omterm{def:g11:functional-groups:alcohol}{alcohol} het koolstofatoom van de \ce{C=O}-groep aan.
Leg met \omterm{def:g11:polarity-and-cohesion:polar-bond}{partiële ladingen} uit waarom dat koolstofatoom een \omterm{def:g12:curly-arrows:site}{acceptorplaats}
is en dat zuurstofatoom een \omterm{def:g12:curly-arrows:site}{donorplaats}. Welk klein \omterm{def:g7:atoms-and-molecules:molecule}{molecuul} gaat er in
totaal af, en tot welke soort behoort de hele reactie?
\end{exercise}

\begin{exercise}[$\star\star\star$]\label{exo:g12:curly-arrows:15}
Etheen reageert met broom \ce{Br2}, een \omterm{def:g11:polarity-and-cohesion:polar-molecule}{apolair molecuul}. Als een
\ce{Br2}-molecuul de \omterm{def:g10:lewis-and-shape:multiple-bond}{dubbele binding} nadert, worden de \omterm{def:g9:inside-the-atom:nucleus}{elektronen} van
\ce{Br-Br} weggeduwd van het dichtstbijzijnde broomatoom. Leg uit waarom
daardoor een \omterm{def:g12:curly-arrows:site}{acceptorplaats} ontstaat, en stel de eerste stap van het
mechanisme voor.
\end{exercise}

\section{Probleem: Broomethaan maken}

\begin{problem}[{Weekendprobleem --- uitgaande van tien gram etheen:
hoeveel broomethaan?}]
\label{pb:g12:curly-arrows:1}
% ledger: aw:C, aw:H, aw:Br, en:H, en:Br, en:C
Broomethaan, een uitgangsstof voor veel \omterm{def:g10:synthesis-yield:synthesis}{syntheses}, wordt gemaakt door
waterstofbromide aan etheen te adderen: \ce{CH2=CH2 + HBr -> CH3-CH2-Br}.
Een scheikundige gebruikt \qty{10.0}{g} etheen en een overmaat
waterstofbromide; het \omterm{def:g10:synthesis-yield:yield}{rendement} is \qty{75}{\percent}.

\medskip
\noindent\textbf{Deel I --- De plaatsen.}
\begin{enumerate}
  \item Teken de \omterm{def:g10:lewis-and-shape:lewis-structure}{lewisstructuren} van etheen en van waterstofbromide.
  \item Waar zit de \omterm{def:g12:curly-arrows:site}{donorplaats} van etheen? Waarom?
  \item Bereken het verschil in \omterm{def:g11:polarity-and-cohesion:electronegativity}{elektronegativiteit} van de
        \ce{H-Br}-binding (H 2.20, Br 2.96). Welk \omterm{def:g7:atoms-and-molecules:atom}{atoom} draagt
        $\delta^+$?
  \item Welk \omterm{def:g7:atoms-and-molecules:atom}{atoom} van \ce{HBr} is de \omterm{def:g12:curly-arrows:site}{acceptorplaats}?
\end{enumerate}

\medskip
\noindent\textbf{Deel II --- Het mechanisme.}
\begin{enumerate}[resume]
  \item Teken de eerste stap met zijn twee \omterm{def:g12:curly-arrows:curly-arrow}{gebogen pijlen}.
  \item Geef de twee deeltjes die ontstaan en hun ladingen.
  \item Controleer het behoud van lading in de eerste stap.
  \item Teken de tweede stap met zijn \omterm{def:g12:curly-arrows:curly-arrow}{gebogen pijl}.
  \item Hoeveel \omterm{def:g9:inside-the-atom:nucleus}{elektronen} omringen het positieve koolstofatoom van het
        intermediair? Waarom is het zo reactief?
\end{enumerate}

\medskip
\noindent\textbf{Deel III --- Wat voor reactie?}
\begin{enumerate}[resume]
  \item Wat is het intermediair van het mechanisme?
  \item Waarom komt het niet voor in de totale \omterm{def:g8:balanced-equations:equation}{reactievergelijking}?
  \item Tot welke soort behoort de reactie?
  \item Verandert de reactie het koolstofskelet, of alleen de
        \omterm{def:g11:functional-groups:functional-group}{karakteristieke groep}?
\end{enumerate}

\medskip
\noindent\textbf{Deel IV --- Hoeveel product?}
\begin{enumerate}[resume]
  \item Bereken de \omterm{def:g10:the-mole:molar-mass}{molaire massa's} van etheen en van broomethaan.
  \item Bereken de hoeveelheid etheen.
  \item Vul de \omterm{def:g10:reaction-progress-table:extent}{voortgangstabel} in. Welke \omterm{def:g7:chemical-reactions:reactant}{beginstof} is beperkend?
  \item Bereken de maximale massa broomethaan.
  \item Bereken de massa die je krijgt bij een \omterm{def:g10:synthesis-yield:yield}{rendement} van
        \qty{75}{\percent}.
  \item Geef het eindantwoord: welke massa broomethaan krijgt de
        scheikundige?
\end{enumerate}
\end{problem}
